\chapter{Cardinals Exponentiation}

While addition and multiplication of cardinals is very simple, since
$$\aleph_\alpha+\aleph_\beta=\aleph_\alpha\cdot\aleph_\beta=\max\{\aleph_\alpha,\aleph_\beta\}$$
exponentiation is way more complex.

In particular, we will see that The Generalized Continuum Hypothesis (GCH):
$$2^{\aleph_\alpha}=\aleph_{\alpha_+1},~\text{ for all }\alpha$$
simplify cardinal exponentiation enormously and, without it, for regular cardinals there isn't much to prove\footnote{in fact, this is a statement in the formal sense. We have a lot of results of the form: giving some relations between cardinals, any possible assignment to those cardinals that doesn't contradicts such relations is realized in some model.} besides
\begin{equation}\label{eq:2_exp}
    2^{\aleph_\alpha}\geq\aleph_{\alpha+1},\qquad\alpha\leq\beta\Rightarrow 2^{\aleph_\alpha}\leq2^{\aleph_\beta}
\end{equation}
and the next result about cofinality.

\section{König's Theorem}

Before proving some theorems in cardinal exponentiation, let's prove one of the most famous results in cardinal arithmetic, known as König's Theorem, which is a strict inequality! Notice how rare this is, the only other strict inequality we saw in cardinal arithmetic was probably Cantor's Theorem. In fact, we will see Cantor is a direct consequence of König's.

\begin{theorem}{König's Theorem}{konig}
    If $\kappa_i<\lambda_i$, where $\kappa_i$ and $\lambda_i$ are cardinal, $\forall i\in I$, then
    $$\sum_{i \in I}\kappa_i<\prod_{i\in I}\lambda_i$$
\end{theorem}

The idea behind the proof is: If $\kappa_i<\lambda_i$, then every function $f_i:\kappa_i\to\lambda_i$ isn't surjective, therefore, we can find an element $b_i\in\lambda_i\backslash\mathrm{Im}(f_i)$, and then we obtain, for each function $f:\bigcup_{i\in I}A_i\to\prod_{i\in I}B_i$, the element $b=\langle b_i:i\in I\rangle$ which is not in the image of $f$.

\begin{proof}
    \todo{should I use daggers?}Let $\langle A_i:i\in I\rangle$ and $\langle B_i:i\in I\rangle$ be such that $|A_i|=\kappa_i$ and $|B_i|=\lambda_i$, with $A_i$ mutually disjoint, $\forall i\in I$. Since $|A_i|<|B_i|$, then, given $f_i:A_i\to B_i$ injective, we know they can't be injective. So, define
    $$B_i^\dagger=\{x\in B_i:x\notin f_i[A_i]\}$$
    Since $f_i$ isn't surjective, $B_i^\dagger\neq\emptyset$, $\forall i\in I$, then, by AC, we can choose $b_i^\dagger\in B_i^\dagger$, $\forall i\in I$. Now, let
    $$F:\bigcup_{i\in I}A_i\to\prod_{i\in I}B_i$$
    be the unique function satisfying
    $$\pi_j(F(a))=\begin{cases}
        f_i(a),~i=j\\
        b_j^\dagger,~i\neq j
    \end{cases}$$
    Now, if $F(x)=F(y)$, with $x\in A_i$ and $y\in A_j$, then
    $$\pi_i(F(x))=\pi_i(F(y))$$
    if $i\neq j$, then
    $$\pi_i(F(x))=f_i(x)=b_i^\dagger=\pi_i(F(y))$$
    a contradiction, since $b_i^\dagger\notin f_i[A_i]$. If $i=j$, then $f_i(x)=f_i(y)$, but, since $f_i$ is injective, $x=y$. This proves $F$ is an injection. Let's now prove there is no surjection.

    Assume by contradiction that there is some $G:\bigcup_{i\in I}A_i\to\prod_{i\in I}B_i$ surjective, and define
    $$B_i^\ddagger=\{x\in B_i:x\notin \pi_i(G[A_i])\}$$
    Since $(\pi_i\circ G):A_i\to B_i$, then it's not surjective and, therefore, $B_i^\ddagger\neq\emptyset$, $\forall i\in I$. By AC there is a $b=\left\langle b_i^\ddagger:i \in I\right\rangle$ such that $b_i^\ddagger\in B_i^\ddagger$. Since $G$ is surjective, there is a $x\in A_i$, for some $i\in I$, such that $G(x)=b^\ddagger$, so $\pi_i(G(x))=\pi_i\rp{b^\ddagger}\in B_i^\ddagger$, a contradiction, then $G$ is not surjective.
\end{proof}

Notice that the choices of $B_i^\dagger$ and $B_i^\ddagger$ in the proof looks like the choice of $X=\{x\in A:x\notin f(x)\}$ in Cantor's Theorem proof and, in fact, Cantor's Theorem is a corollary of König's Theorem for $\kappa_i=1$ and $\lambda_i=2$, $\forall i\in I$.

\subsection{Cardinal Exponentiation for Regular Cardinals}

With such big results in hands, we can prove the following lemma
\begin{lemma}{}{cf_exponentiation}
    For any ordinal $\alpha$
    $$\cf\rp{2^{\aleph_\alpha}}>\aleph_\alpha$$
\end{lemma}
Note that this lemma restricts the possible values $2^{\aleph_\alpha}$ can assume, depending on its cofinality. As an example, which also reinforce the comments in \cref{sec:ramb_cof} and in the last footnote, notice that we can't have $2^{\aleph_0}=\aleph_\omega$, since then we would have $\aleph_0=\cf\rp{\aleph_\omega}=\cf\rp{2^{\aleph_0}}>\aleph_0$, a contradiction.

In fact, Robert Solovay published a result in a really famous one page paper named ``$2^{\aleph_0}$ can be anything it ought to be'' which basically says that the only two things we can prove about the continuum is: that it's uncountable, and that it has uncountable cofinality. Formally, the result says that for any uncountable cardinal $\kappa$ with uncountable cofinality, there is a model of ZFC where this is realized, i.e., where $2^{\aleph_0}=\kappa$.

Nowadays, such result is a direct consequence of a more general one known as Easton's Theorem, that says basically the same thing, but for any regular cardinal, i.e., the only constraints on permissible values for $2^\kappa$ when $\kappa$ is a regular cardinal are \cref{eq:2_exp} and the previous lemma.

\begin{proof}
    Let $\vartheta=\cf\rp{2^{\aleph_\alpha}}$, then, there are $\kappa_\nu<2^{\aleph_\alpha}$, $\nu<\vartheta$ such that
    $$2^{\aleph_\alpha}=\sum_{\nu<\vartheta}\kappa_\nu$$
    Using König's Theorem for $\lambda_\nu=2^{\aleph_\alpha}$, we have
    $$2^{\aleph_\alpha}=\sum_{\nu<\vartheta}\kappa_\nu<\prod_{\nu<\vartheta}2^{\aleph_\alpha}=\rp{2^{\aleph_\alpha}}^\vartheta$$
    Now, assume by contradiction that $\vartheta\leq\aleph_\alpha$, then
    $$2^{\aleph_\alpha}>\rp{2^{\aleph_\alpha}}^\vartheta\leq\rp{2^{\aleph_\alpha}}^{\aleph_\alpha}=2^{\aleph_\alpha\cdot\aleph_\alpha}=2^{\aleph_\alpha}$$
    a contradiction, so $\vartheta=\cf\rp{2^{\aleph_\alpha}}>\aleph_\alpha$.
\end{proof}

Notice how we didn't use anything special about $2$, and then we can repeat the exact same argument for any cardinal $\kappa>1$, giving us:

\begin{corollary}{}{cf_exponentiation}
    For any cardinals $\kappa,\lambda>1$, we have
    $$\cf\rp{\kappa^\lambda}>\lambda$$
\end{corollary}

\subsection{Digression on Exponentiation for Singular Cardinals}

As we saw, the inequalities above are the only ones we can prove about regular cardinals, the next natural step is to ask if there is a generalization, or some similar result about singular cardinals. In fact, for singular cardinals $\kappa$, various additional constraints in the behavior of $2^\kappa$ are known.

By know, the most interesting one is Silver's Theorem, since it implies that Easton's Theorem can't be extended to include singular cardinals. But, nowadays, Shelah's work on PCF Theory gives a lot of results on the possible values of $2^\kappa$, for $\kappa$ singular. It shows that such values are strongly influenced by the values on smaller cardinals, as we will see next and when proving Silver's Theorem, whereas Easton shows this is not the case for regular cardinals.

\begin{lemma}{}{baby_silver}
    Let $\aleph_\alpha$ be a singular cardinal, if $2^{\aleph_\xi}=\aleph_\beta$ for all $\xi<\alpha$, then $2^{\aleph_\alpha}=\aleph_\beta$.
\end{lemma}

\begin{proof}
    Since $\aleph_\alpha$ is singular, there are $\kappa_i<\aleph_\alpha$, $i\in I$, with $|I|=\aleph_\lambda<\aleph_\alpha$ such that
    $$2^{\aleph_\alpha}=\sum_{i\in I}\kappa_i$$
    By hypothesis $2^{\kappa_i}=2^{\aleph_\lambda}=\aleph_\beta$, then
    $$2^{\aleph_\alpha}=2^{\sum_{i\in I}\kappa_i}=\prod_{i\in I}2^{\kappa_i}=\prod_{i\in I}\aleph_\beta=\rp{\aleph_\beta}^{\aleph_\lambda}=\rp{2^{\aleph_\lambda}}^{\aleph_\lambda}=2^{\aleph_\lambda}=\aleph_\beta$$
\end{proof}

\section{Evaluating \texorpdfstring{$\aleph_\alpha^{\aleph_\beta}$}{alepha alephb}}

Let's first state some important lemmas. The first simplify $\aleph_\alpha^{\aleph_\beta}$ for when $\alpha\leq\beta$.

\begin{lemma}{}{card_exp_a_leq_b}
    If $\alpha\leq\beta$, then
    $$\aleph_\alpha^{\aleph_\beta}=2^{\aleph_\beta}$$
\end{lemma}

\begin{proof}
    Since $2<\aleph_\alpha$, we have $2^{\aleph_\beta}\leq\aleph_\alpha^{\aleph_\beta}$. By Cantor's Theorem $\aleph_\alpha<2^{\aleph_\alpha}$, then
    $${\aleph_\alpha}^{\aleph_\beta}\leq\rp{2^{\aleph_\alpha}}^{\aleph_\beta}=2^{\aleph_\alpha\cdot\aleph_\beta}=2^{\aleph_\beta}$$
\end{proof}

The next natural step is trying to calculate $\aleph_\alpha^{\aleph_\beta}$ for when $\alpha>\beta$. The next lemma will show itself useful when trying to do this.

\begin{lemma}{}{card_exp_a_geq_b}
    If $\alpha\geq\beta$, then
    $$\left|[\omega_\alpha]^{\aleph_\beta}\right|=\aleph_\alpha^{\aleph_\beta}$$
\end{lemma}

\begin{proof}
    \todo{review}Since $|\omega_\alpha\times\omega_\beta|=\aleph_\alpha\cdot\aleph_\beta=\aleph_\alpha=|\omega_\alpha|$, then $\left|[\omega_\alpha\times\omega_\beta]^{\aleph_\beta}\right|=\left|[\omega_\alpha]^{\aleph_\beta}\right|$. Now, every function $f:\omega_\beta\to\omega_\alpha$ is in $[\omega_\alpha\times\omega_\beta]^{\aleph_\beta}$, therefore $\omega_\alpha^{\omega_\beta}\subseteq[\omega_\alpha\times\omega_\beta]^{\aleph_\beta}$ and so
    $$\left|\omega_\alpha^{\omega_\beta}\right|=\aleph_\alpha^{\aleph_\beta}\leq\left|[\omega_\alpha\times\omega_\beta]^{\aleph_\beta}\right|=\left|[\omega_\alpha]^{\aleph_\beta}\right|$$
    For the reverse inequality, if $X\in[\omega_\alpha]^{\aleph_\beta}$, then there is some $f$ in $\omega_b$ such that $\mathrm{Im}(f)=X$. With this, given a function $f_X$ for each $X\in[\omega_\alpha]^{\aleph_\beta}$, let $F(X)=f_X$. If $X\neq Y$, then $\mathrm{Im}(f_X)=X$ and $\mathrm{Im}(f_Y)=Y$, so $F(X)=f_X\neq f_Y=F(y)$, i.e., $F$ is injective, then
    $$\left|[\omega_\alpha]^{\aleph_\beta}\right|\leq\left|\omega_\alpha^{\omega_\beta}\right|=\aleph_\alpha^{\aleph_\beta}$$
\end{proof}

\subsection{Under GCH}

Now, as a warmup, let's see what we got under the Generalized Continuum Hypothesis (GCH).

\begin{lemma}{(Under GCH)}{regular_card_exp_gch}
    If $\aleph_\alpha$ is a regular cardinal, then
    $$\aleph_\alpha^{\aleph_\beta}=
    \begin{cases}
        \aleph_\alpha,\text{ if }\alpha>\beta\\
        \aleph_{\beta+1},\text{ if }\alpha\leq\beta
    \end{cases}
    $$
\end{lemma}

\begin{proof}
    \todo{review}If $\alpha\leq\beta$, by \cref{lem:card_exp_a_leq_b} $\aleph_\alpha^{\aleph_\beta}=2^{\aleph_\beta}=\aleph_{\beta+1}$. Let then $\alpha>\beta$. If we let $S=[\omega_\alpha]^{\aleph_\beta}$, by \cref{lem:card_exp_a_leq_b} we have $|S|=\aleph_\alpha^{\aleph_\beta}$. In \cref{sec:weakly_inaccessible}, we saw $\cf(\kappa)$ can be interpreted as the least size of an unbounded set in $\kappa$. In particular, since $\cf(\aleph_\alpha)=\aleph_\alpha$ for regular cardinals, we know every $X\subseteq\omega_\alpha$, with $|X|<\aleph_\alpha$, is limited. Let then
    $$B=\bigcup_{\delta<\omega_\alpha}\mc{P}(\delta)$$
    be the collection of all limited subsets of $\omega_\alpha$. We claim $|B|\leq\aleph_\alpha$; then, since $S\subset B$, we have $|S|=\aleph_\alpha^{\aleph_\beta}=\aleph_\alpha$, since clearly $|S|=\aleph_\alpha^{\aleph_\beta}\geq\aleph_\alpha$. From the definition of $B$ we have
    $$|B|\leq\sum_{\delta<\omega_\alpha}2^{|\delta|}$$
    However, for any cardinal $\aleph_\gamma<\aleph_\alpha$, we also have $2^{\aleph_\gamma}=\aleph_{\gamma+1}\leq\aleph_\alpha$, therefore $2^{|\delta|}\leq\aleph_\alpha$, for any $\delta<\omega_\alpha$, then
    $$|B|\leq\sum_{\delta<\omega_\alpha}2^{|\delta|}\leq\sum_{\delta<\omega\alpha}\aleph_\alpha=\aleph_\alpha\cdot\aleph_\alpha=\aleph_\alpha$$
\end{proof}

Let's now prove a similar, but more complicated, formula for singular cardinals.

\begin{lemma}{(Under GCH)}{singular_card_exp_gch}
    If $\aleph_\alpha$ is a singular cardinal, then
    $$\aleph_\alpha^{\aleph_\beta}=\begin{cases}
        \aleph_\alpha,\text{ if }\aleph_\beta<\cf\rp{\aleph_\alpha},\\
        \aleph_{\alpha+1},\text{ if }\cf\rp{\aleph_\alpha}\leq\aleph_\beta\leq\aleph_\alpha,\\
        \aleph_{\beta+1},\text{ if }\aleph_\beta\geq\aleph_\alpha
    \end{cases}$$
\end{lemma}

Schematically, we obtain \cref{fig:singular_gch}.

\begin{figure}
\centering
\asyinclude{function_graph}
\caption{graph of $f(\aleph_\beta)=\aleph_\alpha^{\aleph_\beta}$ under GCH, for $\aleph_\alpha$ singular.}
\label{fig:singular_gch}
\end{figure}

\begin{proof}
    \todo{review}If $\beta\geq\alpha$, then $\aleph_\beta\geq\aleph_\alpha$ and $\aleph_\alpha^{\aleph_\beta}=2^{\aleph_\beta}=\aleph_{\beta+1}$. If $\aleph_\beta<\cf\rp{\aleph_\alpha}$, then any $X\subseteq\omega_\alpha$ with $|X|=\aleph_\beta<\cf\rp{\aleph_\alpha}<\aleph_\alpha$ is bounded in $\omega_\alpha$ and, therefore, by the same argument as in the previous theorem, we have $\aleph_\alpha^{\aleph_\beta}=\aleph_\alpha$. Let then $\cf\rp{\aleph_\alpha}\leq\aleph_\beta\leq\aleph_\alpha$. On the one hand
    \begin{equation}\label{eq:aux_eq}
    \aleph_\alpha\leq\aleph_\alpha^{\aleph_\beta}\leq\aleph_\alpha^{\aleph_\alpha}=2^{\aleph_\alpha}=\aleph_{\alpha+1}
    \end{equation}
    But on the other, by \cref{cor:cf_exponentiation}, $\cf\rp{\aleph_\alpha^{\aleph_\beta}}>\aleph_\beta$ and, since $\aleph_\beta\geq\cf\rp{\aleph_\alpha}$, then $\cf\rp{\aleph_\alpha^{\aleph_\beta}}\neq\cf\rp{\aleph_\alpha}$, so $\aleph_\alpha^{\aleph_\beta}\neq\aleph_\alpha$. With this and \cref{eq:aux_eq}, the only possible value for $\aleph_\alpha^{\aleph_\beta}$ is $\aleph_{\alpha+1}$.
\end{proof}

\subsection{Hausdorff's Formula}

Lastly, let's see what we can do without GCH.

\begin{theorem}{Hausdorff's Formula}{hausdorff_formula}
    For any ordinals $\alpha$ and $\beta$, we have
    $$\aleph_{\alpha+1}^{\aleph_\beta}=\aleph_\alpha^{\aleph_\beta}\cdot\aleph_{\alpha+1}$$
\end{theorem}

\begin{proof}
    \todo{review}If $\beta\geq\alpha+1$, then, since $\aleph_{\alpha+1}\leq\aleph_\beta\leq2^{\aleph_\beta}$, we have
    $$\aleph_{\alpha+1}^{\aleph_\beta}=2^{\aleph_\beta}=2^{\aleph_\beta}\cdot\aleph_{\alpha+1}=\aleph_\alpha^{\aleph_\beta}\cdot\aleph_{\alpha+1}$$
    Let then $\beta\leq\alpha$. Since $\aleph_\alpha^{\aleph_\beta}\leq\aleph_{\alpha+1}^{\aleph_\beta}$ and $\aleph_{\alpha+1}\leq\aleph_{\alpha+1}^{\aleph_\beta}$, then
    $$\aleph_\alpha^{\aleph_\beta}\cdot\aleph_{\alpha+1}\leq\aleph_{\alpha+1}^{\aleph_\beta}\cdot\aleph_{\alpha+1}^{\aleph_\beta}=\aleph_{\alpha+1}^{\aleph_\beta}$$
    Therefore, it's sufficient to show $\aleph_{\alpha+1}^{\aleph_\beta}\leq\aleph_\alpha^{\aleph_\beta}\cdot\aleph_{\alpha+1}$. Since $\omega_\beta<\omega_{\alpha+1}$ and $\omega_{\alpha+1}$ is regular, then for any $f:\omega_\beta\to\omega_{\alpha+1}$, $f$ is bounded, i.e., there is some $\gamma<\omega_{\alpha+1}$ such that $f(\xi)<\gamma$, for every $\xi<\omega_\beta$. With this
    $$\omega_{\alpha+1}^{\omega_\beta}=\bigcup_{\gamma<\omega_{\alpha+1}}\gamma^{\omega_\beta}$$
    Now, any $\gamma<\omega_{\alpha+1}$ is such that $|\gamma|\leq\aleph_\alpha$ and, since
    $$\left|\bigcup_{\gamma<\omega_{\alpha+1}}\gamma^{\omega_\beta}\right|\leq\sum_{\gamma<\omega_{\alpha+1}}|\gamma|^{\aleph_\beta}$$
    then
    $$\aleph_{\alpha+1}^{\aleph_\beta}\leq\sum_{\gamma<\omega_{\alpha+1}}|\gamma|^{\aleph_\beta}\leq\sum_{\gamma<\omega_{\alpha+1}}\aleph_\alpha^{\aleph_\beta}=\aleph_{\alpha}^{\aleph_\beta}\cdot\aleph_{\alpha+1}$$
\end{proof}
